Probleme/Fragen:

(1) Marginalien sind bei mehreren Wörtern teilweise sehr gesperrt, z.B. S. 2 Ein-, zwei-, dreiwertige Verben; S. 23 Kasus und Präposition; 

(2) Etwas aus dem Fließtext ragt in den Rand (Signal schwarzes Kästchen): bei automatischer Trennung ragt der Trennstrich in den Rand (schwarzes Kästchen) z.B. S. 8: Grund-modell; auch ohne Trennung S. 30; 31 (zweimal); S. 33 bei Trennung am Bindestrich Nominativ-NP; S. 43, S. 45 bei Bindestrich; S. 47 bei Marginalie; S. 49 Tabelle ragt in den Rand – Problem wiederholt sich 

(3) S. 130 Tabelle 1. Spalte: getrenntes "Präpositional-objekt": Problem mit Zeilenabstand; gleiches Problem S. 233 in Tabelle 1., 2., 4. Spalte "Subjekts-prädikativ"

(4) Tabelle S. 238 zu groß? – schwarzer Streifen an der Seite 
 
(5) S. 451: Satzpunkt steht in einer separaten Zeile.

(6) Können leere Felder in den Analysetabellen (z.B. S. 137, 150, 157) grau hinterlegt/schraffiert werden? 

(7) Farbe: In Abschnitt 8.5.1. sind MRT-Aufnahmen in Farbe. Der Text ist so gestaltet, dass man die Abbildungen auch in einen SW-Ausdruck verstehen könnte. Dürfen die Abbildungen in Farbe bleiben, oder müssen sie nach SW umgewandelt werden? 


